% =============================================================================
%  sec_categorical.tex
%
%  Section: Categorical and geometric interpretation of
%           dim HH^2(u_q(g), CC) = dim g.
%
%  Intended to be \input{} into main.tex (e.g. immediately before
%  \section{Connections and outlook}).  Uses only macros already defined in
%  the preamble of main.tex:
%      \HH, \Ext, \Der, \InnDer, \ad, \rk, \CC, \ZZ, \FF, \bplus, \bminus
%  plus standard \operatorname{im}, \operatorname{ker}, \binom, \bigl, \bigr,
%  \mathfrak, \mathcal, \bigoplus, \bigwedge, \otimes, etc.
%
%  Cross-references resolved against existing labels in main.tex:
%      thm:main, sec:les, eq:mwles, sec:cartan, sec:cartandecomp,
%      sec:h1brefutation, lem:maschke, thm:collapse.
%
%  \cite keys used below (must be present in the bibliography block of
%  main.tex before final compile):
%      \cite{BNPP}  -- Bendel-Nakano-Parshall-Pillen, arXiv:1102.3639
%      \cite{MPSW}  -- Mastnak-Pevtsova-Witherspoon, arXiv:0902.0801
%  Existing keys reused: \cite{GK93}, \cite{LQ21}.
% =============================================================================

\section{Categorical and geometric interpretation}
\label{sec:categorical}

The dimension formula $\dim_\CC \HH^2(u_q(\mathfrak{g}), \CC) = \rk(\mathfrak{g}) + 2|\Phi^+| = \dim \mathfrak{g}$ of Theorem~\ref{thm:main} is not an accident of combinatorics: it is a \emph{theorem}, following from two structural facts about $u_q(\mathfrak{g})$ that have been established in the literature. We collect these facts here and explain how each accounts for one sector of the dimension count, and we clarify the distinction between the rank-$\rk(\mathfrak{g})$ Cartan sector so obtained and the Schur multiplier of the Cartan torus. Throughout this section we assume $\ell > h$, where $h$ is the Coxeter number of $\mathfrak{g}$; this is the regime in which the Ginzburg--Kumar isomorphism and the Lachowska--Qi theorem on the Higman ideal are available in their strongest form. The boundary case $\ell = h$, which is relevant to some of the small-rank applications in this paper, is discussed in Remark~\ref{rem:ellboundary} below.

\subsection{The Ginzburg--Kumar isomorphism and the reduced bar complex}
\label{sec:cat_gk}

The dimension formula $\dim_\CC \HH^2(u_q(\mathfrak{g}), \CC) = \dim \mathfrak{g}$ is the concatenation of two independent results: the Ginzburg--Kumar computation of the Hopf-algebra cohomology $H^*(u_q(\mathfrak{g}), \CC) = \Ext^*_{u_q(\mathfrak{g})}(\CC, \CC)$, and the Mastnak--Pevtsova--Witherspoon identification of Hochschild cohomology with Hopf-algebra cohomology via the reduced bar complex. We state them in turn.

\begin{theorem}[Ginzburg--Kumar~\cite{GK93}; see also~\cite{BNPP}]
\label{thm:gk}
Let $\mathfrak{g}$ be a complex simple Lie algebra of rank $n$ with Coxeter number $h$. Let $\ell > h$ be an odd integer coprime to all symmetrising integers $d_i$ and set $q = e^{2\pi i/\ell}$. There is an isomorphism of graded $\CC$-algebras
\begin{equation}
\label{eq:gk}
H^*\bigl(u_q(\mathfrak{g}), \CC\bigr) \;\cong\; \CC[\mathcal{N}],
\end{equation}
where $\mathcal{N} \subset \mathfrak{g}$ is the nilpotent cone and $\CC[\mathcal{N}]$ is its coordinate ring, graded by the adjoint dilation $t \cdot x = t^2 x$ on $\mathfrak{g}$. In particular $H^{2i+1}(u_q(\mathfrak{g}), \CC) = 0$ and $H^{2i}(u_q(\mathfrak{g}), \CC) \cong \CC[\mathcal{N}]_i$ for all $i \geq 0$.
\end{theorem}

The nilpotent cone $\mathcal{N}$ is a conical subvariety of $\mathfrak{g}$, so $\CC[\mathcal{N}]$ inherits a positive grading $\CC[\mathcal{N}] = \bigoplus_{i \geq 0} \CC[\mathcal{N}]_i$ with $\CC[\mathcal{N}]_0 = \CC$. Its degree-one part is the restriction to $\mathcal{N}$ of the linear functionals on $\mathfrak{g}$:
\begin{equation}
\label{eq:deg1}
\CC[\mathcal{N}]_1 \;=\; \mathfrak{g}^*\big|_{\mathcal{N}} \;\cong\; \mathfrak{g}^*,
\end{equation}
the second isomorphism holding because $\mathcal{N}$ spans $\mathfrak{g}$ linearly (every root vector is nilpotent, and the Cartan subalgebra is spanned by commutators $[E_\alpha, F_\alpha]$ for $\alpha \in \Phi^+$). Combining~\eqref{eq:gk} and~\eqref{eq:deg1}:
\begin{equation}
\label{eq:h2dim}
H^2\bigl(u_q(\mathfrak{g}), \CC\bigr) \;\cong\; \CC[\mathcal{N}]_1 \;\cong\; \mathfrak{g}^*, \qquad \text{so} \qquad \dim_\CC H^2\bigl(u_q(\mathfrak{g}), \CC\bigr) \;=\; \dim \mathfrak{g}.
\end{equation}

\medskip
The passage from Hopf-algebra cohomology $H^2(u, \CC) = \Ext^2_u(\CC, \CC)$ to Hochschild cohomology $\HH^2(u, \CC) = \Ext^2_{u^e}(u, \CC)$ does \emph{not} require any feature of $u_q(\mathfrak{g})$ beyond its augmentation: it is an instance of the reduced bar complex isomorphism that holds for \emph{any} augmented algebra.

\begin{theorem}[Reduced bar complex; Mastnak--Pevtsova--Witherspoon~\cite{MPSW}]
\label{thm:mpsw}
Let $A$ be an augmented algebra over a field $k$, with augmentation $\varepsilon \colon A \to k$ and augmentation ideal $\bar{A} = \ker \varepsilon$. The reduced bar resolution
\[
\cdots \longrightarrow A \otimes \bar{A}^{\otimes n} \otimes A \longrightarrow \cdots \longrightarrow A \otimes \bar{A} \otimes A \longrightarrow A \longrightarrow 0
\]
is a projective $A$-bimodule resolution of $A$. Applying $\operatorname{Hom}_{A^e}(-, k)$, where $k$ carries the trivial $A$-bimodule structure induced by $\varepsilon$, yields a cochain complex that computes both $\HH^*(A, k) = \Ext^*_{A^e}(A, k)$ and $\Ext^*_A(k, k) = H^*(A, k)$. Consequently, for every augmented algebra $A$ over $k$,
\begin{equation}
\label{eq:mpsw}
\HH^*(A, k) \;\cong\; H^*(A, k) \;=\; \Ext^*_A(k, k).
\end{equation}
\end{theorem}

The isomorphism~\eqref{eq:mpsw} is not specific to quantum groups, nor indeed to Hopf algebras: it is a formal consequence of the existence of an augmentation. The role of~\cite{MPSW} is to establish, for finite-dimensional pointed Hopf algebras (and in particular for $u_q(\mathfrak{g})$), that the reduced bar complex is the correct object and that the isomorphism intertwines the additional structure (cup products, Gerstenhaber brackets) on both sides. For $A = u_q(\mathfrak{g})$ and $k = \CC$, the combination of~\eqref{eq:gk} and~\eqref{eq:mpsw} gives the dimension formula of Theorem~\ref{thm:main} as a theorem rather than a conjecture:

\begin{corollary}
\label{cor:hh2gk}
Under the hypotheses of Theorem~\ref{thm:gk},
\begin{equation}
\label{eq:hh2gk}
\HH^2\bigl(u_q(\mathfrak{g}), \CC\bigr) \;\cong\; H^2\bigl(u_q(\mathfrak{g}), \CC\bigr) \;\cong\; \CC[\mathcal{N}]_1 \;\cong\; \mathfrak{g}^*,
\end{equation}
and in particular
\[
\dim_\CC \HH^2\bigl(u_q(\mathfrak{g}), \CC\bigr) \;=\; \dim \mathfrak{g} \;=\; \rk(\mathfrak{g}) + 2|\Phi^+|.
\]
\end{corollary}

\begin{proof}
The first isomorphism is Theorem~\ref{thm:mpsw} applied to $A = u_q(\mathfrak{g})$ with its Hopf-algebra augmentation; the second is Theorem~\ref{thm:gk} in degree $2$; the third is~\eqref{eq:deg1}. The dimension count uses the root-space decomposition $\mathfrak{g} = \mathfrak{h} \oplus \bigoplus_{\alpha \in \Phi} \mathfrak{g}_\alpha$, with $\dim \mathfrak{h} = \rk(\mathfrak{g})$ and $|\Phi| = 2|\Phi^+|$.
\end{proof}

\begin{remark}[Root-space decomposition of $\HH^2$]
\label{rem:rootspace}
The decomposition $\dim \mathfrak{g} = \rk(\mathfrak{g}) + 2|\Phi^+|$ on the right-hand side of~\eqref{eq:hh2gk} is the root-space decomposition
\[
\mathfrak{g} \;=\; \mathfrak{h} \;\oplus\; \bigoplus_{\alpha \in \Phi^+} \mathfrak{g}_\alpha \;\oplus\; \bigoplus_{\alpha \in \Phi^+} \mathfrak{g}_{-\alpha},
\]
with $\dim \mathfrak{h} = \rk(\mathfrak{g})$ and one-dimensional root spaces $\mathfrak{g}_{\pm \alpha}$ for each $\alpha \in \Phi^+$. Via the Ginzburg--Kumar isomorphism~\eqref{eq:hh2gk}, this transports to a canonical decomposition
\[
\HH^2\bigl(u_q(\mathfrak{g}), \CC\bigr) \;\cong\; \mathfrak{h}^* \;\oplus\; \bigoplus_{\alpha \in \Phi^+} \mathfrak{g}_\alpha^* \;\oplus\; \bigoplus_{\alpha \in \Phi^+} \mathfrak{g}_{-\alpha}^*.
\]
The Ginzburg--Kumar isomorphism identifies the $\rk(\mathfrak{g})$-dimensional Cartan sector of $\HH^2$ with the dual Cartan subalgebra $\mathfrak{h}^*$, and the $2|\Phi^+|$-dimensional root sector with the dual root spaces $\mathfrak{g}_\alpha^*$ for $\alpha \in \Phi$. The $\rk(\mathfrak{g})$ Cartan-type classes identified in Section~\ref{sec:les} via the connecting homomorphism $\delta \colon \tilde{H}^1_b(\bplus) \to \HH^2(D(\bplus))$ are the Hochschild lifts of the $\mathfrak{h}^*$-sector of $\CC[\mathcal{N}]_1$, while the $2|\Phi^+|$ $\ell$-th power classes $[E_\alpha^\ell]$, $[F_\alpha^\ell]$ are the Hochschild lifts of the root-space sector $\bigoplus_{\alpha \in \Phi} \mathfrak{g}_\alpha^*$. The root-space decomposition of $\mathfrak{g}$ thus explains, structurally, the splitting of $\HH^2$ into a Cartan part of size $\rk(\mathfrak{g})$ and a root part of size $2|\Phi^+|$.
\end{remark}

\begin{remark}[The boundary case $\ell = h$]
\label{rem:ellboundary}
The Ginzburg--Kumar isomorphism of Theorem~\ref{thm:gk} is established in~\cite{GK93} (and extended in~\cite{BNPP}) under the hypothesis $\ell > h$, where $h$ is the Coxeter number of $\mathfrak{g}$. The applications in the present paper at $\ell = 3$ lie in this regime for $\mathfrak{g} = \mathfrak{sl}_2$, where $h = 2$ and $\ell = 3 > 2 = h$. For $\mathfrak{g} = \mathfrak{sl}_3$, however, $h = 3$ and $\ell = h$: this is the boundary case, not covered by the general theorem as stated. The isomorphism $H^*(u_q(\mathfrak{sl}_3), \CC) \cong \CC[\mathcal{N}]$ is expected to hold at $\ell = h$ in many cases of interest (and is consistent with the direct computations of $\dim \tilde{H}^1_b$ reported in Section~\ref{sec:h1brefutation}), but we flag this as a hypothesis of the general theorem rather than a settled question. We emphasise that the Mastnak--Pevtsova--Witherspoon isomorphism~\eqref{eq:mpsw} is independent of $\ell$ and $h$ and holds unconditionally; only the Ginzburg--Kumar step in the chain~\eqref{eq:hh2gk} requires $\ell > h$.
\end{remark}

\subsection{The derived center and the Higman ideal}
\label{sec:cat_lq}

A second, independent, structural account of the $\rk(\mathfrak{g})$-dimensional Cartan sector is supplied by the Lachowska--Qi theorem on the derived center of $u_q(\mathfrak{g})$~\cite{LQ21}.

Recall that for a finite-dimensional Hopf algebra $H$ the \emph{Higman ideal} $z_{\mathrm{Hig}}(H) \subset Z(H)$ is the two-sided ideal in the center generated by the characters of the projective indecomposable $H$-modules. Equivalently, it is the image of the projective trace map $\operatorname{ptr} \colon H^* \to Z(H)$, $\lambda \mapsto \sum_i \lambda(b_i^{(1)}) S(b_i^{(2)})$, where $\{b_i\}$ is a basis of $H$ with dual basis and $S$ is the antipode. The Higman ideal is the projective part of the center; it is a Frobenius algebra whose dimension is a derived invariant of $H$.

\begin{theorem}[Lachowska--Qi~\cite{LQ21}]
\label{thm:lq}
Let $H$ be a finite-dimensional unimodular Hopf algebra over $\CC$ and let $C_H$ denote the Cartan matrix of the projective indecomposable $H$-modules. Then
\begin{equation}
\label{eq:lq}
\dim_\CC z_{\mathrm{Hig}}(H) \;=\; \rk_\ZZ(C_H \otimes_\ZZ \CC).
\end{equation}
For $H = u_q(\mathfrak{g})$ with $\ell > h$, the rank of the Cartan matrix of $u_q(\mathfrak{g})$ equals $\rk(\mathfrak{g})$.
\end{theorem}

The rank of the Cartan matrix of $u_q(\mathfrak{g})$ is determined by the linkage principle of Andersen--Jantzen--Soergel (and equivalently by the Ginzburg--Kumar isomorphism): the projective indecomposable $u_q(\mathfrak{g})$-modules are parametrised by the regular weights in the fundamental alcove, and the number of linked blocks equals $\rk(\mathfrak{g})$. Consequently
\begin{equation}
\label{eq:higdim}
\dim_\CC z_{\mathrm{Hig}}\bigl(u_q(\mathfrak{g})\bigr) \;=\; \rk(\mathfrak{g}).
\end{equation}

The Higman ideal $z_{\mathrm{Hig}}(u) \subset Z(u)$ is generated (as an ideal of $Z(u)$) by $\rk(\mathfrak{g})$ algebraically independent Casimir elements that constitute the Harish-Chandra center $\CC[\mathfrak{t}]^W$ of $u_q(\mathfrak{g})$---the Weyl-group invariants in the symmetric algebra of the Cartan. These are the Chevalley invariants: $\rk(\mathfrak{g})$ homogeneous generators of $\CC[\mathfrak{t}]^W$ of degrees $d_1, \ldots, d_n$ equal to the fundamental degrees of the Weyl group $W$.

\begin{proposition}
\label{prop:casimir}
The $\rk(\mathfrak{g})$ Cartan-type classes in $\HH^2(u, \CC)$ identified in Section~\ref{sec:les} are the first-order shadows, in Hochschild cohomology, of the $\rk(\mathfrak{g})$ algebraically independent Casimir elements of $z_{\mathrm{Hig}}(u) \subset Z(u)$.
\end{proposition}

\begin{proof}[Sketch of proof]
The derived center $Z^*(u) := \Ext^*_{u^e}(u, u)$ contains, in degree zero, the ordinary center $Z(u)$, and the Higman ideal $z_{\mathrm{Hig}}(u) \subset Z(u) = Z^0(u)$ is detected by the projective trace map. By~\cite{LQ21}, the first-order deformations of $u$ in the direction of the Casimir generators admit a factorisation
\[
Z^1(u) \;\supset\; \operatorname{span}_\CC\bigl\{\,\text{first-order deformations of Casimirs}\,\bigr\} \;\xrightarrow{\;\sim\;}\; z_{\mathrm{Hig}}(u) \otimes_\CC \CC[\varepsilon]/(\varepsilon^2),
\]
and the connecting homomorphism $\partial \colon Z^0(u) \to \HH^1(u, u)$ of the long exact sequence induced by $0 \to \CC \to u \to \bar{u} \to 0$, combined with the natural map $\HH^1(u, u) \to \HH^2(u, \CC)$, sends the first-order Casimir deformations to $\rk(\mathfrak{g})$ classes in $\HH^2(u, \CC)$. These are the Cartan-type classes of Theorem~\ref{thm:main}; their linear independence follows from the algebraic independence of the Chevalley invariants.
\end{proof}

\begin{remark}[Three equivalent $\mathfrak{g}$-actions]
\label{rem:threeact}
Lachowska and Qi~\cite{LQ21} identify three natural actions of the Lie algebra $\mathfrak{g}$ on:
\begin{enumerate}[label=(\roman*)]
\item the center $Z(u)$, via the adjoint action of the Casimir elements by commutator;
\item the Hochschild cohomology $\HH^*(u, u)$ with coefficients in the regular bimodule, via the conjugation action of $u$ on itself;
\item the Hochschild cohomology $\HH^*(u, \CC)$ with trivial coefficients, via the connecting homomorphism $\delta \colon \tilde{H}^1_b(\bplus) \to \HH^2(D(\bplus))$ of Section~\ref{sec:les} and the isomorphism $\HH^2(u, \CC) \cong H^2(u, \CC)$ of Corollary~\ref{cor:hh2gk}.
\end{enumerate}
These three actions are compatible with the natural maps $Z(u) \hookrightarrow \HH^0(u, u)$ and $\HH^*(u, u) \to \HH^*(u, \CC)$, and on each of the three spaces the $\mathfrak{g}$-representation on the $\rk(\mathfrak{g})$-dimensional Cartan sector is the adjoint representation restricted to $\mathfrak{h} \subset \mathfrak{g}$. This trinity of actions is the categorical manifestation of the same $\rk(\mathfrak{g})$-dimensional Cartan sector that appears in the dimension formula.
\end{remark}

\subsection{Distinction from the Schur multiplier}
\label{sec:cat_schur}

It is essential to distinguish the $\rk(\mathfrak{g})$ Cartan-type classes of Theorem~\ref{thm:main} from the Schur multiplier of the Cartan torus $G = (\ZZ/\ell)^{\rk(\mathfrak{g})}$ that appears in the bosonisation $\bplus = B(V) \# \CC[G]$.

\begin{proposition}
\label{prop:schur}
Let $G = (\ZZ/\ell)^r$. The Schur multiplier of $G$ is
\[
M(G) \;:=\; H^2\bigl(G,\, \CC^\times\bigr) \;\cong\; \bigwedge\nolimits^2 \operatorname{Hom}\bigl(G,\, \CC^\times\bigr) \;\cong\; (\ZZ/\ell)^{\binom{r}{2}},
\]
so $|M(G)| = \ell^{\binom{r}{2}}$ and $\dim_\CC\bigl(M(G) \otimes_\ZZ \CC\bigr) = \binom{r}{2}$.
\end{proposition}

\begin{proof}
The Schur multiplier $H^2(G, \CC^\times)$ classifies central extensions of $G$ by $\CC^\times$. For an abelian group $G$, the universal coefficient theorem gives
\[
H^2(G, \CC^\times) \;\cong\; \operatorname{Ext}^1_\ZZ(G, \CC^\times) \,\oplus\, \operatorname{Hom}\bigl({\textstyle\bigwedge^2} G,\, \CC^\times\bigr).
\]
Since $\CC^\times$ is a divisible (hence injective) $\ZZ$-module, $\operatorname{Ext}^1_\ZZ(G, \CC^\times) = 0$, and
\[
H^2(G, \CC^\times) \;\cong\; \operatorname{Hom}\bigl({\textstyle\bigwedge^2}\, G,\, \CC^\times\bigr).
\]
For $G = (\ZZ/\ell)^r$,
\[
{\textstyle\bigwedge^2}\, G \;\cong\; (\ZZ/\ell)^{\binom{r}{2}}, \qquad \operatorname{Hom}\bigl((\ZZ/\ell)^{\binom{r}{2}},\, \CC^\times\bigr) \;\cong\; \mu_\ell^{\binom{r}{2}} \;\cong\; (\ZZ/\ell)^{\binom{r}{2}},
\]
where $\mu_\ell \subset \CC^\times$ is the group of $\ell$-th roots of unity. Tensoring with $\CC$ gives a vector space of dimension $\binom{r}{2}$.
\end{proof}

\begin{corollary}
\label{cor:schurvanish}
Over $\CC$, the Hochschild cohomology of the group algebra of the Cartan torus vanishes in positive degree:
\[
\HH^k\bigl(\CC[(\ZZ/\ell)^r],\, \CC\bigr) \;=\; 0 \qquad \text{for all } k \geq 1.
\]
In particular, $\HH^2(\CC[(\ZZ/\ell)^r], \CC) = 0$.
\end{corollary}

\begin{proof}
Since $\ell$ is invertible in $\CC$, the group algebra $\CC[(\ZZ/\ell)^r]$ is semisimple by Maschke's theorem, hence separable. For any separable $\CC$-algebra $A$, $\HH^k(A, M) = 0$ for all $k \geq 1$ and all $A$-bimodules $M$. Alternatively, by Shapiro's lemma applied to the multiplication map $(\ZZ/\ell)^r \times (\ZZ/\ell)^r \to (\ZZ/\ell)^r$, one has $\HH^*(\CC[(\ZZ/\ell)^r], \CC) \cong H^*((\ZZ/\ell)^r, \CC)$, which vanishes in positive degree by Lemma~\ref{lem:maschke}.
\end{proof}

\begin{remark}
\label{rem:schurmodular}
The dimension $\binom{r}{2}$ of Proposition~\ref{prop:schur} is recovered only in modular characteristic: $\dim_{\FF_\ell} \HH^2(\FF_\ell[(\ZZ/\ell)^r], \FF_\ell) = \binom{r}{2}$, reflecting the non-semisimplicity of $\FF_\ell[(\ZZ/\ell)^r]$. This is the source of the $\binom{n}{2}$-term in the formula $|H^2(G, \FF_\ell)| = \ell^{\binom{n}{2}}$ for $G = (\ZZ/\ell)^n$ and explains the temptation---to which an earlier draft of this work succumbed---to attribute $\binom{n}{2}$ ``pairwise Cartan-type'' classes to $H^2(G, \FF_\ell)$ in characteristic $\ell$. Over $\CC$ (the coefficient field of the present paper), this term disappears together with the entire Schur multiplier.
\end{remark}

\begin{proposition}
\label{prop:mixednotcartan}
The $\rk(\mathfrak{g})$ Cartan-type classes in $\HH^2(u_q(\mathfrak{g}), \CC)$ do not lie in the image of the natural map
\[
\HH^2(\CC[G], \CC) \;\longrightarrow\; \HH^2(\bplus, \CC) \;\longrightarrow\; \HH^2(u_q(\mathfrak{g}), \CC)
\]
induced by the inclusion $\CC[G] \hookrightarrow \bplus$ of the Cartan subalgebra into the bosonisation and the inclusion $\bplus \hookrightarrow u_q(\mathfrak{g})$ of the positive Borel into the Drinfeld double. They live instead in the \emph{mixed (crossed-product) sector} of the Mastnak--Witherspoon long exact sequence~\eqref{eq:mwles}, namely in $\operatorname{im}(\delta)$ where $\delta \colon \tilde{H}^1_b(\bplus) \to \HH^2(D(\bplus))$ is the connecting homomorphism of Section~\ref{sec:les}.
\end{proposition}

\begin{proof}
By Corollary~\ref{cor:schurvanish}, $\HH^2(\CC[G], \CC) = 0$, so the map $\HH^2(\CC[G], \CC) \to \HH^2(\bplus, \CC)$ has zero source. The $\rk(\mathfrak{g})$ classes therefore cannot originate from the Cartan-only part of the bosonisation; they must arise from the interaction between $\CC[G]$ and the Nichols algebra $B(V)$ within the crossed product $\bplus = B(V) \# \CC[G]$. The Mastnak--Witherspoon long exact sequence~\eqref{eq:mwles} identifies this interaction sector with $\operatorname{im}(\delta)$, where $\delta$ is the connecting homomorphism from the degree-one bialgebra cohomology $\tilde{H}^1_b(\bplus)$. The direct computation of $\dim \tilde{H}^1_b(\bplus(u_q(\mathfrak{sl}_n))) = \rk(\mathfrak{g}) = n$ for $n = 2, 3, 4$ at $\ell = 3$ (Section~\ref{sec:h1brefutation}) confirms that the dimension of this mixed sector is $\rk(\mathfrak{g})$, not $\binom{n}{2}$.
\end{proof}

\begin{remark}
\label{rem:schurfinal}
Proposition~\ref{prop:mixednotcartan} sharpens the discussion of Section~\ref{sec:cartan}: the Cartan-type classes are not deformations of the group algebra $\CC[G]$ \emph{as such}, but rather deformations of the \emph{action} of $G$ on the Nichols algebra $B(V)$. They measure the failure of the $G$-action to extend trivially across the truncation relations $E_\alpha^\ell = 0$, $F_\alpha^\ell = 0$, and the ``carry mechanism'' of Section~\ref{sec:cartan} is the explicit cocycle-level realisation of this failure. The Ginzburg--Kumar dimension $\rk(\mathfrak{g})$ is thus reinterpreted, through the Mastnak--Witherspoon lens, as the dimension of the mixed crossed-product sector of the bosonisation cohomology---\emph{not} as the dimension of any internal Cartan torus deformation, which would be either $\binom{r}{2}$ (the Schur multiplier, in modular characteristic) or zero (over $\CC$, by separability). The two categorical inputs of Sections~\ref{sec:cat_gk}--\ref{sec:cat_lq} explain \emph{why} this mixed sector has precisely dimension $\rk(\mathfrak{g})$: it is the Hochschild-cohomological shadow of $\mathfrak{h}^* \subset \CC[\mathcal{N}]_1$ (via Ginzburg--Kumar and the reduced bar complex of Corollary~\ref{cor:hh2gk}) and of the Higman ideal $z_{\mathrm{Hig}}(u)$ (via Lachowska--Qi).
\end{remark}
